<%= heading2 %>{UML diagrams}

UML diagrams can be written as text next to the \LaTeX{} source and are typeset in the fonts of the document.
\href{https://ctan.org/pkg/metauml}{MetaUML} draws class, component, and further diagrams.
It is MetaPost code, which \texttt{lualatex} runs itself via the \href{https://ctan.org/pkg/luamplib}{luamplib} package; no further tools are needed.
Thus, the MetaUML diagrams are shown only when compiling with \texttt{lualatex}.
The \href{https://ctan.org/pkg/pgf-umlsd}{pgf-umlsd} package draws sequence diagrams with TikZ and works with every engine.
A comparison of the packages for UML diagrams is given at \url{https://github.com/latextemplates/gadr-latex-packages/blob/main/uml-diagrams.md}.

\ifluatex
  <%- bexample %>
  \begin{figure}
    \centering
    \sffamily
    \begin{mplibcode}
      input metauml;
      beginfig(1);
      Class.Driver("Driver")()();
      Class.Car("Car")("+model: String")("+drive()");
      Class.Wheel("Wheel")()();
      Class.Person("Person")()();
      % distances between the boxes in PostScript points (1/72 inch)
      leftToRight(70)(Driver, Car, Wheel);
      topToBottom(50)(Car, Person);
      drawObjects(Driver, Car, Wheel, Person);
      link(association)(Driver.e -- Car.w);
      item(iAssoc)("drives >")(obj.s = .5[Driver.e, Car.w]);
      % the diamond is drawn at the end of the path
      link(composition)(Wheel.w -- Car.e);
      item(iAssoc)("have 4 >")(obj.s = .5[Car.e, Wheel.w]);
      link(association)(Car.s -- Person.n);
      item(iAssoc)("< owns")(obj.w = .5[Car.s, Person.n] + (5, 0));
      endfig;
    \end{mplibcode}
    \caption{Class diagram drawn with MetaUML}
    \label{fig:uml-class}
  \end{figure}
  <%- eexample %>

  <%- bexample %>
  \begin{figure}
    \centering
    \sffamily
    \begin{mplibcode}
      input metauml;
      beginfig(1);
      Component.Client("Client")();
      Component.Server("Server")();
      % MetaUML has no automatic layout: each object is placed relative to another one;
      % here, the left edge of Server is 140 PostScript points (1/72 inch, about 5 cm) right of the right edge of Client
      Server.w = Client.e + (140, 0);
      drawObjects(Client, Server);
      % the interface: Server provides it (ball), Client requires it (socket)
      pair api;
      api = .5[Client.e, Server.w];
      link(providedInterface)(Server.w -- api);
      link(requiredInterface)(Client.e -- api);
      item(iAssoc)("API")(obj.s = api + (0, 8));
      endfig;
    \end{mplibcode}
    \caption{Component diagram drawn with MetaUML}
    \label{fig:uml-component}
  \end{figure}
  <%- eexample %>
\fi

<%- bexample %>
\begin{figure}
  \centering
  \sffamily
  \begin{sequencediagram}
    \newthread{client}{Client}
    % [3]: horizontal gap (in cm) to the box of the previous instance; the default (0.2) leaves no room for the message labels
    \newinst[3]{server}{Server}
    \begin{call}{client}{request()}{server}{response}
    \end{call}
  \end{sequencediagram}
  \caption{Sequence diagram drawn with pgf-umlsd}
  \label{fig:uml-sequence}
\end{figure}
<%- eexample %>
